\begin{figure}[H]
    \centering
    \begin{tikzpicture}[x=0.62cm, y=1.0cm]

        \node[blkreg, minimum width=1.7cm, minimum height=0.52cm] (heap) at (-2.3,0.95)
            {\scriptsize σωρός $k$};
        \node[notemute, font=\tiny, anchor=north, align=center] at (-2.3,0.62)
            {ρίζα $=\tau$};

        \foreach \i in {0,...,7} { \node[lanereg] (t\i) at (\i,0.95) {$\tau$}; }
        \node[notemute, font=\tiny, anchor=west] at (8.35,0.95) {κατώφλι σε καταχωρητή};
        \draw[flow] (heap.east) -- (t0.west);

        \foreach \b/\y/\vals/\mask in {%
            1/-0.15/{3,1,4,2,5,1,3,2}/{0,0,0,0,0,0,0,0},
            2/-0.80/{2,4,1,3,2,5,1,4}/{0,0,0,0,0,0,0,0},
            3/-1.45/{1,3,2,4,1,2,3,5}/{0,0,0,0,0,0,0,0},
            4/-2.10/{2,1,9,3,1,4,2,1}/{0,0,1,0,0,0,0,0}} {
            \foreach \val [count=\i from 0] in \vals {
                \node[lane] (b\b-\i) at (\i,\y) {$\val$};
            }
            \node[rowlab, font=\tiny] at (b\b-0.west) {$B_\b$};
            \foreach \mv [count=\i from 0] in \mask {
                \node[font=\scriptsize, minimum width=0.62cm] at (\i+9.4,\y)
                    {\ifnum\mv=1\relax\textcolor{figdrmD}{\bfseries 1}\else\textcolor{figink}{0}\fi};
            }
        }

        \node[notemute, font=\tiny, anchor=south] at (12.9,-0.02) {μάσκες};

        \draw[bracemirror] (9.1,-2.42) -- (16.7,-2.42);
        \node[note, font=\scriptsize, anchor=north] at (12.9,-2.62)
            {ένωση τεσσάρων μασκών};

        \node[blkdrop, minimum width=4.4cm, minimum height=0.46cm, anchor=east] (skip) at (7.6,-3.62)
            {\scriptsize $=0$: προσπέρασμα και των τεσσάρων};
        \node[blkcac, minimum width=5.2cm, minimum height=0.46cm, anchor=west] (sift) at (9.3,-3.62)
            {\scriptsize $\neq 0$: κατολίσθηση μόνο για τα αναμμένα};

        \draw[flow] (12.9,-3.00) -- (4.05,-3.00) -- (4.05,-3.39);
        \draw[flow] (12.9,-3.00) -- (13.49,-3.00) -- (13.49,-3.39);

        \node[note, anchor=north, align=center] at (7.0,-4.15)
            {Στη συνηθισμένη περίπτωση ακολουθείται σχεδόν πάντοτε η αριστερή διαδρομή.};

    \end{tikzpicture}
    \caption{Η μαζική απόρριψη του \tl{map-reduce}. Η ρίζα του σωρού, δηλαδή το
    χειρότερο από τα \tl{k} στοιχεία που έχουν ήδη επιλεγεί, αντιγράφεται σε
    ολόκληρο έναν διανυσματικό καταχωρητή και συγκρίνεται ταυτόχρονα με τέσσερα
    μπλοκ της εισόδου, παράγοντας μία δυαδική μάσκα ανά μπλοκ. Αν η ένωση των
    τεσσάρων μασκών είναι μηδενική, και τα τέσσερα μπλοκ προσπερνιούνται με έναν
    μόνο έλεγχο. Βαθμωτή εργασία εκτελείται αποκλειστικά για τα \tl{bits} που είναι
    αναμμένα. Όταν ο σωρός έχει ήδη γεμίσει με μεγάλες τιμές, το μεγαλύτερο μέρος
    της εισόδου διασχίζεται χωρίς κανένα στοιχείο να εξεταστεί ποτέ ξεχωριστά.}
    \label{fig:cpu-mapreduce-reject}
\end{figure}
